\section{Introduction}\label{sec-intro}

\subsection{Context}
An independent price verification (IPV) function controls the parameters that front office (FO) uses to value
and risk-manage equity derivatives. For repo, the reference is not observed directly. It is \emph{implied}: Totem
publishes consensus prices of calls and puts, from which forwards are extracted by put--call parity, and a
separate Totem service publishes consensus dividends. The repo curve that reconciles the two is then compared with
FO's repo marks, and the differences are multiplied by FO sensitivities to produce control lines (variances,
reserves, prudent-valuation adjustments).

Two features of Totem make this delicate. The option service and the dividend service are fed by different
contributors, and within each service the number of contributors changes with tenor; a typical single stock might
have fifteen contributors at one year on options and ten on dividends, fewer beyond two years. The implied repo
mixes the two consensus means, and so inherits the noise of both and the artefacts of their mismatch. Controllers
see this as irregular shapes and unjustified jumps in implied repo curves.

\subsection{Questions answered}
This note answers four questions:
\begin{enumerate}[label=(\alph*)]
\item \textbf{Where do irregular implied repo curves come from?} \Cref{sec-anatomy} decomposes the implied repo
exactly into the option panel's repo plus a panel gap, and derives the signature of each mechanism (panel gap,
contributor drop-off, ex-date timing, stale ex-date profile, dividend bucket error, spot snapshot, gross/net
definition).
\item \textbf{Can a better statistical method recover the ``true'' repo?} No: \cref{thm-nonid} shows that
consensus data cannot identify the split between repo and dividends of the option panel.
\item \textbf{Do we need a clean Totem repo curve to control repo?} For forward-linear exposure, no:
\cref{thm-control,thm-exact} prove that the total of the dividend and repo lines does not depend on the dividend
panel, and \cref{prop-materiality} measures, from sensitivities alone, how much the split matters for a given book.
\item \textbf{How should bad curves be detected?} \Cref{sec-noise,sec-detection} build a structural detector:
a noise model computed from Totem's own statistics, one matched-filter template per mechanism, a calibrated
false-alarm rate and an explanation in natural units for each alarm. \Cref{sec-benchmark} benchmarks it, and
\cref{sec-alternatives} explains where B-spline coefficients, Gaussian mixtures and small neural networks fit.
\end{enumerate}

\subsection{Main results at a glance}
\begin{insight}
\begin{itemize}[leftmargin=12pt]
\item To first order, $b_{\text{implied}} = \bA + (\DA - \DB)/\bigl(T(S-D)\bigr)$ (\cref{prop-gap}). The gap alone
gives \vGapOneY\,bp of repo noise at one year and \vGapFiveY\,bp at five years for a typical single stock.
\item The mechanisms have distinct signatures in \emph{carry space} $B=b\,T$: a constant (spot), a step (dividend
amount, panel change), a spike (forward outlier, ex-date timing), a seasonal pattern (allocation profile).
\item For forward-linear books, the full-revaluation Totem P\&L is \emph{exactly} independent of the dividend panel
when expiries sit on Totem pillars (\cref{thm-exact}); the line decomposition reproduces this to first order
(\cref{thm-control}). The dividend panel only matters through the non-forward exposure $g = s_D + s_b J$.
\item The detector's test statistics are exact matched filters (\cref{thm-glrt}). One variance scale per category,
calibrated on clean history, sets the false-alarm rate exactly (\cref{prop-kappa}).
\item In the benchmark (\vNTest{} synthetic test names, every detector at \vFPR{} false alarms), the structural
detector and a Gaussian mixture on B-spline coefficients of the Totem--FO difference detect a similar share of
anomalies on average (\vMeanStruct{} and \vMeanGmmDiff). The structural detector also finds forward
outliers (\vFwdStruct{} against at most \vFwdShape), explains \vAttr{} of its alarms correctly, and does not mistake
genuine FO mismarks for bad Totem curves (\vFoStruct{} against \vFoGmmDiff).
\end{itemize}
\end{insight}

\subsection{Reading guide}
\Cref{sec-model} fixes notation and proves the elementary facts the rest relies on (bootstrap, parity,
allocation). \Cref{sec-anatomy} is the diagnosis; \cref{sec-control} the control identity; \cref{sec-noise,sec-detection}
the detector; \cref{sec-benchmark} the evidence; \cref{sec-alternatives} the comparison with other approaches;
\cref{sec-ops} the implementation. Proofs follow each statement; longer derivations are in \cref{app-proofs}.
Every numerical example is produced by the code in the repository (\cref{app-repro}).
